% ============================================================================
% Fichier  : partie6/P06_C19_audit_certification.tex
% Titre    : L'Audit du SMSI et la Certification
% Auteur   : MINKA MI NGUIDJOÏ Thierry Emmanuel
% ============================================================================
% Sources : ISO 19011:2018 ; ISO/IEC 27006:2015 ; ISO/IEC 27001 clause 9.2-9.3.
% Mairie d'Eséka : cas fictif.
% ============================================================================

\chapter{L'Audit du SMSI et la Certification}
\label{chap:audit}
\chaptermark{Audit et certification}

\epigraph{%
    « L'audit ne cherche pas le coupable~:\\
    il cherche l'écart entre ce qui est dit et ce qui est fait. »%
}{Principe du cours}

\niveauChapitre{Intermédiaire}{Le Chapitre~\ref{chap:indicateurs}.}

\begin{boxobjectifs}
\begin{itemize}
    \item Distinguer audit interne et audit de certification.
    \item Décrire la méthodologie d'audit (ISO 19011).
    \item Caractériser constat, preuve, non-conformité et écart.
    \item Expliquer la revue de direction.
    \item Décrire le cycle de certification \norme{ISO/IEC 27001}.
\end{itemize}
\end{boxobjectifs}

\bigskip

% ============================================================================
\section*{Question directrice : comment faire reconnaître que la sécurité est maîtrisée ?}
% ============================================================================

La mairie pense son \sigle{SMSI} solide. Un audit le vérifie, et la
certification l'atteste. Ce chapitre présente l'audit, la revue de direction
et le chemin vers la certification.

\medskip
\textbf{Plan du chapitre.}
\begin{enumerate}
    \item Les types d'audit (section~\ref{sec:c19-types}).
    \item La méthodologie (section~\ref{sec:c19-methodo}).
    \item Constat, preuve, écart (section~\ref{sec:c19-constat}).
    \item La revue de direction (section~\ref{sec:c19-revue}).
    \item Le cycle de certification (section~\ref{sec:c19-cycle}).
\end{enumerate}

\separateur

% ============================================================================
\section{Les Types d'Audit}
\label{sec:c19-types}
% ============================================================================

\begin{tableaucours}{L{3.2cm}L{9.0cm}}{%
    \tableauHeader{Type} &
    \tableauHeader{Qui et pourquoi}
}
Audit interne & Mené par l'organisation, exigé par la clause~9.2 \\
Audit de seconde partie & Mené par un client ou une tutelle \\
Audit de certification & Mené par un organisme accrédité (tierce partie) \\
\end{tableaucours}
\vspace{-0.8em}
\begin{center}{\small\itshape Tableau~19.1. Les types d'audit}\end{center}

L'\termecle{audit interne} prépare~: il détecte les écarts avant l'auditeur
externe. L'\termecle{audit de certification} atteste, auprès des tiers, que le
\sigle{SMSI} satisfait la norme.

% ============================================================================
\section{La Méthodologie}
\label{sec:c19-methodo}
% ============================================================================

La norme \norme{ISO 19011} guide l'audit des systèmes de management~:
planification, revue documentaire, entretiens et observations sur site,
rédaction du rapport, suivi des actions. L'auditeur est \termecle{indépendant}
de ce qu'il audite et procède par \termecle{échantillonnage}.

\begin{boxprocedure}
\textbf{Constatation.} Pour l'opération \textsc{Matomb}, le cabinet
\textsc{Matomb Conseil} conduit l'audit interne de la mairie avant la
certification.

\medskip
\textbf{Analyse.} L'audit interne est une répétition générale~: il cherche les
écarts pendant qu'il est encore temps de les corriger. Un écart trouvé en
interne se corrige~; le même écart trouvé en certification coûte un report.
\end{boxprocedure}

% ============================================================================
\section{Constat, Preuve, Écart}
\label{sec:c19-constat}
% ============================================================================

\begin{boxdef}
\termecle{Preuve d'audit}~: enregistrement vérifiable (document, observation,
déclaration). \termecle{Constat}~: résultat de l'évaluation des preuves face à
un critère. \termecle{Non-conformité}~: écart entre l'exigence et la réalité,
\emph{majeure} (défaillance du système) ou \emph{mineure} (manquement isolé).
\end{boxdef}

L'auditeur sépare le \emph{fait} (ce qu'il constate, preuve à l'appui) de son
\emph{appréciation}. Cette rigueur, qui distingue la constatation de l'analyse,
est au cœur du métier d'auditeur~\cite{minka2026auditologie}.

\begin{boxmetier}{Auditeur SMSI}
\textbf{Pas de constat sans preuve.} L'auditeur ne note jamais une impression~:
il adosse chaque constat à une preuve (un document, un relevé, une observation).
Un écart sans preuve n'est pas un écart~: c'est un avis.
\end{boxmetier}

% ============================================================================
\section{La Revue de Direction}
\label{sec:c19-revue}
% ============================================================================

\begin{boxdef}
\termecle{Revue de direction} (clause~9.3)~: examen périodique du \sigle{SMSI}
par la direction, qui passe en revue les résultats des audits, des mesures et
des incidents, et décide des améliorations et des ressources.
\end{boxdef}

La revue de direction est le moment où la direction \emph{reprend la main}~:
elle transforme les constats en décisions. Sans elle, l'audit resterait sans
suite. À Eséka, le maire la préside deux fois par an.

% ============================================================================
\section{Le Cycle de Certification}
\label{sec:c19-cycle}
% ============================================================================

\begin{tableaucours}{L{3.6cm}L{8.6cm}}{%
    \tableauHeader{Étape} &
    \tableauHeader{Contenu}
}
Audit initial (étape 1) & Revue documentaire, préparation \\
Audit initial (étape 2) & Audit sur site, décision de certification \\
Audits de surveillance & Chaque année, maintien du certificat \\
Renouvellement & Au bout de trois ans, nouvel audit complet \\
\end{tableaucours}
\vspace{-0.8em}
\begin{center}{\small\itshape Tableau~19.2. Le cycle de certification ISO/IEC 27001}\end{center}

La certification se \emph{gagne} en deux étapes et se \emph{maintient} par des
audits de surveillance. Elle n'est pas acquise une fois pour toutes~: le
\sigle{SMSI} doit continuer de tourner (Chapitre~\ref{chap:pdca}).

\separateur

\begin{boxresume}
\textbf{Ce qu'il faut retenir de ce chapitre~:}
\begin{itemize}
    \item L'\textbf{audit interne} prépare~; l'\textbf{audit de certification}
          atteste.
    \item La \textbf{méthodologie} (ISO 19011) procède par preuves,
          échantillonnage et indépendance.
    \item Un \textbf{constat} s'adosse à une \textbf{preuve}~; une
          \textbf{non-conformité} est majeure ou mineure.
    \item La \textbf{revue de direction} transforme les constats en décisions.
    \item La \textbf{certification} se gagne en deux étapes et se maintient par
          surveillance.
\end{itemize}
\end{boxresume}

\begin{boxinfo}
\textbf{Points de contrôle de l'élève.} Je sais~: distinguer les types
d'audit~; décrire la méthodologie~; définir preuve, constat, non-conformité~;
expliquer la revue de direction~; décrire le cycle de certification.
\end{boxinfo}

\bigskip

\begin{boxexo}[title={\ding{45}\quad Exercice~19.1. Types d'audit \niveaubase}]
Distinguez audit interne et audit de certification.
\end{boxexo}

\begin{boxexo}[title={\ding{45}\quad Exercice~19.2. Preuve \niveaubase}]
Qu'est-ce qu'une preuve d'audit~? Donnez deux exemples.
\end{boxexo}

\begin{boxexo}[title={\ding{45}\quad Exercice~19.3. Le cycle \niveaubase}]
Énoncez les étapes du cycle de certification \norme{ISO/IEC 27001}.
\end{boxexo}

\begin{boxexo}[title={\ding{45}\quad Exercice~19.4. Majeure ou mineure \niveaumoyen}]
Distinguez une non-conformité majeure d'une mineure, par un exemple à la
mairie.
\end{boxexo}

\begin{boxexo}[title={\ding{45}\quad Exercice~19.5. Répétition générale \niveaumoyen}]
Pourquoi l'audit interne est-il une « répétition générale »~?
\end{boxexo}

\begin{boxexo}[title={\ding{45}\quad Exercice~19.6. Constat sans preuve \niveaumoyen}]
Pourquoi un écart sans preuve n'est-il qu'un avis~?
\end{boxexo}

\begin{boxexo}[title={\ding{45}\quad Exercice~19.7. Indépendance \niveaumoyen}]
Pourquoi l'auditeur doit-il être indépendant de ce qu'il audite~?
\end{boxexo}

\begin{boxexo}[title={\ding{45}\quad Exercice~19.8. Revue de direction \niveaumoyen}]
Que décide la direction lors de la revue, et à partir de quoi~?
\end{boxexo}

\begin{boxexo}[title={\ding{45}\quad Exercice~19.9. Plan d'audit \niveauexpert}]
Rédigez le plan d'audit interne du \sigle{SMSI} de la mairie d'Eséka (périmètre,
critères, étapes).
\end{boxexo}

\begin{boxexo}[title={\ding{45}\quad Exercice~19.10. Fiche d'écart \niveauexpert}]
Rédigez une fiche de non-conformité (constat, preuve, exigence, gravité, action
corrective) pour un écart de la mairie.
\end{boxexo}

\bigskip

\begin{boxinfo}
\textbf{Pour aller plus loin.}
\begin{itemize}
    \item ISO 19011~\cite{iso19011}~; ISO/IEC 27006~\cite{iso27006}~;
          auditologie~\cite{minka2026auditologie}.
\end{itemize}
\end{boxinfo}

\begin{boxinfo}
\textbf{Transition.} Reste le dernier maillon du cycle~: apprendre des
incidents. Le Chapitre~\ref{chap:incidents} clôt le cours.
\end{boxinfo}
